\documentclass[11pt]{article}
\usepackage{icmlko}
\begin{document}

\title{규칙이 죽인 것을 세어 보니 한 건이었다 --- 그리고 그것이 내가 통과시키고 싶은 건이라서, 규칙을 안 고쳤다}
\author{월드모델 실험실}
\date{노트 542 · 2026년 8월 4일}
\maketitle

\begin{abstract}
노트 541 이 노트 446 거부권을 대체할 자(\emph{같은 해 쌍에서도 버는가})를
세우고 알려진 날짜 구멍으로 검증했다. 규칙을 바꾸기 전에 \textbf{그 규칙이
죽인 것을 전수 조사}했다.

\begin{center}
\begin{tabular}{lr}
\hline
 & \\
\hline
원장에서 $446$ 을 언급한 항목 & $14$ 개 \\
\textbf{거부권이 실제로 죽인 것} & $\mathbf{2}$ 건 \\
그중 이미 되살아난 것(노트 $514$) & $1$ 건 \\
\textbf{지금 죽어 있는 것} & $\mathbf{1}$ \textbf{건} \\
\hline
\end{tabular}
\end{center}

\textbf{그 한 건이 내가 통과시키고 싶은 건이다.} 그래서 \textbf{규칙을 안
고친다.}
\end{abstract}

\begin{figure}[t]
\centering
\includegraphics[width=\textwidth]{figs/f1.pdf}
\caption{\textbf{가} --- 전수 조사. \textbf{나} --- 셋 중 둘만 섰다.
\textbf{다} --- 대신 한 것.}
\end{figure}

\section{노트 511 은 이미 살아 있었다}

$446$ 으로 죽은 두 건 중 하나는 \textbf{노트 511 의 ``줄이는 쪽 가중''}이다
--- 판 $+0.0006$ $\cdot$ 날짜 통제 $-0.0005$ $\cdot$ \textbf{KR 만화
$+0.0203$ ($40/40$, 단일 변경에서 본 제일 센 전이 신호)}. 그때 보류됐다.

\begin{quote}
그런데 \textbf{노트 513 이 배선 결함을 찾았다} --- \texttt{DOMDROP} 이 예보
전용으로 걸려 있었다. \textbf{고친 배선에서 다시 재니 날짜 통제가
$-0.0005 \to -0.0000$($18/40$)로 안 움직였고}, 거부권 전제가 사라져
\textbf{노트 514 에서 채택}됐다. 지금 챔피언에 들어 있다.
\end{quote}

\textbf{그러니 $446$ 이 지금 죽이고 있는 것은 시장팝업 \texttt{category}
하나뿐이다}(노트 524 의 \texttt{TIMEAX} 는 세 자 모두 음수라 $446$ 없이도
기각된다).

\section{셋 중 둘만 섰다}

규칙을 바꿀 수 있는 조건을 셋으로 적으면

\begin{center}
\begin{tabular}{ll}
\hline
조건 & \\
\hline
새 자가 양성 대조로 검증됐나 & \textbf{예}(노트 541) \\
규칙이 틀렸다는 증거가 있나 & \textbf{예}(연도 $p = 0.168$ $\cdot$ 하락은 전부 이웃) \\
\textbf{소급 범위가 하나보다 많나} & \textbf{아니오 --- $1$ 건} \\
\hline
\end{tabular}
\end{center}

\begin{quote}
\textbf{소급 범위가 한 건이고 그것이 지금 내가 원하는 건이면, 규칙 변경은
그 자체로 결과를 고르는 것이다.} 노트 512 가 같은 자리에서 적었다 ---
``조항을 결과를 보고 고치지 않는다. 대신 조항이 겨냥한 것을 더 정확히
겨냥하는 실험을 설계한다.''
\end{quote}

\section{대신 한 것 --- 여섯째 자와 재검토 문턱}

\begin{itemize}
\item \textbf{``같은 해 이득''을 관문의 여섯째 자로 넣는다.} 앞으로 모든
후보가 이 값을 함께 적는다. 거부권 $446$ 은 \textbf{그대로 둔다} --- 새 자는
보고되지만 거부권을 이기지 못한다.
\item \textbf{재검토 문턱을 미리 적는다.} \textbf{독립된 세 건}에서
``$446$ 이 우는데 같은 해 이득이 $0$ 밖 양수''가 나오면 그때 변경을 사전
등록한다. \textbf{지금은 $1/3$.}
\end{itemize}

\textbf{시장팝업 \texttt{category} 는 기각 유지}다. 자료는 살아 있다 ---
그 도메인 $\rho$ 를 $+0.157$ 에서 $+0.315$ 로 만드는 한 칸이고, \emph{붙이는
자리}를 더 찾는다(노트 540 이 공유 열과 전용 열 둘을 닫았다).

\section*{조항}

\begin{quote}
\textbf{규칙을 바꾸려면 그 규칙이 죽인 것이 여럿이어야 한다.} 소급 범위가
한 건이고 그것이 지금 원하는 건이면, 규칙 변경은 결과를 고르는 것이다.
\textbf{대신 새 자를 추가하고 재검토 문턱을 미리 적는다.}
\end{quote}

\section*{남은 것}

\begin{itemize}
\item \textbf{제삼의 형태} --- \texttt{grp} 도 전용 열도 아닌 자리. 규칙을
안 건드리고 다섯 자를 통과하는 길이 있는지(진행 중).
\item \textbf{팝업 $\cdot$ 아이돌 $\cdot$ 만화도 \texttt{grp} 덮음이 $0\%$}
인데, \texttt{grpaxes.SPEC} 에는 만화가 들어 있다. 왜 $0\%$ 인지 안 봤다.
\item \textbf{$1/3$ 이 $3/3$ 이 되는 데 얼마나 걸리나} --- 그때까지 이 자는
보고만 된다.
\end{itemize}

\end{document}
